\providecommand{\ConstructiveWindingCertificateUp}{\mathsf{ConstructiveWindingCertificate}^{\uparrow}}

\chapter{A Constructive Winding-Certificate Layer for the Argument-Principle Route}
\label{ch:visions-rh-constructive-winding-certificate-layer}
\origin{ai}

The argument-principle route at finite height separates into two obligations.
The topological and combinatorial half defines a winding number for rational
boundary polygons and proves that this number is well defined under the
certificate operations used by a finite verifier. The analytic half proves
that such a polygon faithfully tracks $\zeta(\partial R)$ and that the
constructive argument-principle theorem identifies the resulting winding with
the number of zeros inside $R$. Only the topological half is closed here.

The present layer therefore does not prove the Riemann Hypothesis, does not
locate every zeta zero, and does not supply the analytic bridge from
$\zeta(\partial R)$ to a rational polygon. It records the completed
mathlib-free, 0-axiom, propext-free Lean certificate layer in the namespace
$\mathsf{BEDC.Derived.RHRoute.RationalPolygonWinding}$ and its sibling
modules. The layer works with finite rational vertices and the signed crossing
count of the positive real ray. Its `Up' versions use the BEDC integer carrier
so that the algebra remains inside the same constructive kernel fragment.

\begin{definition}[Rational-polygon winding certificate layer]
\label{def:rh-rational-polygon-winding-certificate-layer}
A rational-polygon winding certificate consists of a finite cyclic list of
rational complex vertices, the oriented edge list obtained from that cycle,
and the signed sum of edge crossings of the positive real ray. The ordinary
integer count is written as $\mathsf{computedWinding}$, while the BEDC-integer
count is written as $\mathsf{computedWindingUp}$. A certificate generator may
perturb rational vertices within certified crossing margins, refine an edge by
inserting an intermediate rational point, rotate the starting vertex, reverse
edge orientations, and split a rectangular boundary into adjacent boxes.
\end{definition}

\begin{principle}[Perturbation stability]
\label{prin:rh-winding-perturbation-stability}
If every edge of a rational polygon carries a margin certificate keeping it
away from the crossing discriminants, then any edgewise rational perturbation
preserving those margins has the same computed winding. This is the stability
fact used when a finite evaluator replaces a boundary trace by nearby rational
vertices without changing the topological certificate.
\end{principle}
\leanchecked{BEDC.Derived.RHRoute.RationalPolygonWindingStability.computedWinding\_stable\_of\_edgewiseMarginPreserving}

\begin{principle}[Collinear edge subdivision]
\label{prin:rh-winding-collinear-edge-subdivision}
Inserting a rational point strictly between the endpoints of a segment splits
the segment into two collinear directed segments without changing the total
crossing contribution of that edge. This is the local algebraic fact behind
boundary refinement.
\end{principle}
\leanchecked{BEDC.Derived.RHRoute.RationalPolygonEdgeSubdivision.edgeCrossing\_collinear\_split}

\begin{principle}[Surface-level subdivision invariance]
\label{prin:rh-winding-subdivision-invariance}
At the summed BEDC-integer level, replacing one edge in an oriented edge list
by its two collinear subedges preserves the total crossing sum. Thus adaptive
refinement changes the mesh of the polygon, not its winding certificate.
\end{principle}
\leanchecked{BEDC.Derived.RHRoute.RationalPolygonEdgeSubdivision.sumEdgeCrossingsUp\_subdivide\_one}

\begin{principle}[Cyclic rotation invariance]
\label{prin:rh-winding-cyclic-rotation}
The chosen starting vertex of a nonempty cyclic vertex list is not part of the
certificate content. Rotating the list by one step preserves the BEDC-integer
winding count.
\end{principle}
\leanchecked{BEDC.Derived.RHRoute.RationalPolygonWinding.computedWindingUp\_cyclic\_rotate\_one}

\begin{principle}[Orientation reversal cancellation]
\label{prin:rh-winding-orientation-cancellation}
Traversing the same rational edge in the two opposite directions contributes
opposite signed crossings, so the two directed copies have net contribution
$0$ in the BEDC integer carrier.
\end{principle}
\leanchecked{BEDC.Derived.RHRoute.RationalPolygonWinding.edgeCrossingUp\_reverse\_cancel}

\begin{principle}[Box-split additivity]
\label{prin:rh-winding-box-split-additivity}
When a rectangular boundary is split into adjacent boxes, the shared interior
edge is traversed once in each direction and cancels. The resulting crossing
sum is additive across the split, giving the combinatorial skeleton for
recursive box certificates.
\end{principle}
\leanchecked{BEDC.Derived.RHRoute.RationalPolygonWinding.sumEdgeCrossingsUp\_box\_split}

Together these checked facts say the following precise thing: for rational
boundary polygons, the winding number is a finite combinatorial object that is
well defined, stable under certified perturbation, invariant under edge
subdivision and cyclic rotation, antisymmetric under orientation reversal, and
additive under box splitting. These are the structural properties needed by
the topological side of an argument-principle certificate, and they are all
kernel-verified in the 0-axiom BEDC fragment. They are not a zero-counting
theorem for $\zeta$.

\begin{principle}[Open analytic bridges]
\label{prin:rh-winding-open-analytic-bridges}
Two analytic obligations remain open. First, the route needs a located
interval evaluator for $\zeta$ on $\partial R$. The intended evaluator passes
through the eta series
$$
\eta(s)=\sum_{n\geq1}(-1)^{n-1}n^{-s},
\qquad
\zeta(s)=\frac{\eta(s)}{1-2^{1-s}},
$$
with
$$
n^{-s}=\exp(-s\ln n),
\qquad
n^{-it}=\cos(t\ln n)-i\sin(t\ln n).
$$
The critical point is the tail: a naive Leibniz estimate does not directly
apply because the phase $n^{-it}$ rotates in the complex plane. A strict
Euler--Boole alternating-summation bound is the concrete 0-axiom shape needed
for the critical strip, including derivative bounds for interval-Newton or
Krawczyk certificates.

Second, the route needs the constructive argument-principle theorem itself:
a discrete Rouch\'e or faithful-approximation theorem proving that a rational
polygon whose certified tube contains $\zeta(\partial R)$ has the same winding
as the true boundary trace and therefore counts zeros inside $R$. These are
multi-week constructive-analysis obligations, not completed certificate
payload.
\end{principle}

A usable Euler--Boole obligation can be stated without hiding the analytic
work. For a box
$$
B=\{s=\sigma+it:\sigma_-\leq\sigma\leq\sigma_+,\ t_-\leq t\leq t_+\},
\qquad 0<\sigma_-<\sigma_+<1,
$$
set $X=N+1$ and $f_s(x)=x^{-s}$. With Euler polynomials defined by
$$
\frac{2e^{xz}}{e^z+1}
=
\sum_{\nu\geq0}E_\nu(x)\frac{z^\nu}{\nu!},
$$
Euler--Boole summation gives endpoint terms plus a remainder
$$
R_N(s)
=
(-1)^N
\left[
\frac12\sum_{\nu=0}^{2L-1}
\frac{E_\nu(0)}{\nu!}\,
f_s^{(\nu)}(X)
+\mathcal E_L(s)
\right],
$$
where
$$
f_s^{(\nu)}(X)=(-1)^\nu(s)_\nu X^{-s-\nu},
\qquad
E_{2j}(0)=0\quad(j\geq1).
$$
A constructive interval evaluator must turn this into a uniform complex-box
bound for $\mathcal E_L(B)$ and for $\partial_s\mathcal E_L(B)$, then propagate
those enclosures through $\zeta(s)=\eta(s)/(1-2^{1-s})$. This is the analytic
bottleneck shared by the winding and local-zero routes.

\begin{principle}[Located-zero complement]
\label{prin:rh-winding-located-zero-complement}
The local alternative is to locate a single zero by a Krawczyk or
interval-Newton certificate. Near the first nontrivial zero,
$c\approx 1/2+14.1347i$, one chooses a high-precision approximate inverse
$a\approx 1/\zeta'(c)$ and studies the Newton-like map
$$
N(z)=z-a\,\zeta(z).
$$
For a closed complex box $B$ around $c$, the constructive target has the form
$$
N(B)\subseteq B,
\qquad
\sup_{z\in B}|1-a\zeta'(z)|<1.
$$
This traps a unique zero in $B$. It is a local located-zero certificate, not a
global statement about all zeros. It is complementary to the rational-polygon
winding layer because both legs need the same located-$\zeta$ interval
evaluator and the same Euler--Boole complex tail control.
\end{principle}

This should also be read against the route separation in
\autoref{ch:visions-rh-dyadic-truncation-lp-route-refutation}. The sharp
dyadic Laguerre--P\'olya premise fails as a global preservation strategy. The
present chapter does not try to revive that premise. It keeps the finite
topological certificate layer and points the analytic work toward located
$\zeta$ evaluation, Krawczyk-style local boxes, and a constructive
argument-principle bridge.

\begin{closurestatus}{\ConstructiveWindingCertificateUp}
  \constructivestory{A $\ConstructiveWindingCertificateUp$ packet is generated from finite rational boundary vertices, oriented rational edges, signed positive-ray crossing rows, perturbation-margin rows, subdivision rows, cyclic-rotation rows, reversed-edge cancellation rows, and box-split rows, with all exported structure witnessed by the listed Lean targets.}
  \theoryclosure{\scopedClosure}
  \scopeclosed{Rational-boundary-polygon winding number's combinatorial certificate layer, including well-definedness, stability, subdivision and cyclic-rotation invariance, orientation cancellation, and box additivity; excluding located-zeta evaluation and the analytic argument-principle theorem}
  \formalstatus{\theoremCheckedV}
  \leantarget{BEDC.Derived.RHRoute.RationalPolygonWinding.sumEdgeCrossingsUp\_box\_split}
  \bridgestatus{none}
  \notclaimed{Not closed: the located-zeta interval evaluator and the constructive argument-principle theorem identifying winding with zero count}
  \upgradepath{Add located complex exp/ln/trig continuity evaluators plus discrete Rouch\'e or faithful-approximation error control}
  \origin{ai}
\end{closurestatus}
